% !PRINTED 676
\weberpage{676}

%<<CONT
in $S$ enthalten. Denn gibt es ein nicht in $S$ enthaltenes Polygon
$\mathfrak{A}_{s+1}$ in $A$, so ist

\[
(\mathfrak{A}_1, \mathfrak{A}_2, \ldots, \mathfrak{A}_s, \mathfrak{A}_{s+1})
\]

eine Schar von der Dimension $s+1$. Damit ist bewiesen:

7. \gesperrt{Die Polygone einer Klasse bilden eine Schar
von einer endlichen Dimension $s$. Diese Zahl soll
die Dimension der Klasse heißen.}

Wenn es in einer Klasse $C$ Polygone gibt, die durch ein
Polygon $\mathfrak{A}$ einer Klasse $A$ teilbar sind, so ist $C$ durch $A$ teilbar.
Ist

\[
C = AB,
\]

so erhält man die Dimension von $B$ dadurch, daß man die Schar
aus $C$ aussucht, die durch irgend ein Polygon $\mathfrak{A}$ in $A$ teilbar
ist. Diese Dimension ist also nur von den beiden Klassen $A$
und $C$ abhängig und soll mit $(A, C)$ bezeichnet werden. Das
Zeichen für die Dimension einer Klasse $A$ ist hiernach $(O, A)$,
wenn $O$ die Klasse des Nullecks ist. Wir haben dann

\weq{14}{(A, C) = (A, AB) = (O, B),}

und wenn $a$ die Ordnung der Klasse $A$ ist, so ist nach 5.:

\weq{15}{(A, C) \geqq (O, C) - a.}

Bezeichnen wir mit $A$ die eigentliche Klasse, die man erhält,
wenn man $\mathfrak{M}$ überall weghebt, so soll die uneigentliche Klasse
mit dem Teiler $\mathfrak{M}$ durch $\mathfrak{M}A$ bezeichnet sein.

8. \gesperrt{Der Teiler einer uneigentlichen Klasse ist ein
isoliertes Polygon.}

Ist nämlich $\mathfrak{M}A$ eine uneigentliche Klasse, so läßt sich in
$A$ ein Polygon $\mathfrak{A}$ finden, das relativ prim zu $\mathfrak{M}$ ist. Ist $\mathfrak{M}'$ mit
$\mathfrak{M}$ äquivalent, so ist $\mathfrak{M}'\mathfrak{A}$ in $\mathfrak{M}A$ enthalten und mithin durch
$\mathfrak{M}$ teilbar. Es ist also auch $\mathfrak{M}'$ durch $\mathfrak{M}$ teilbar, also mit $\mathfrak{M}$
identisch. Es gibt also in der Klasse von $\mathfrak{M}$ nur das einzige
Polygon $\mathfrak{M}$.

\begin{center}
§ 193. \gesperrt{Normalbasen.}
\end{center}

Nach dem Vorigen können wir eine Funktion $z$ des Körpers
$\Omega$ von der $n$ten Ordnung in der Form

\weq{1}{z = \frac{\mathfrak{U}'}{\mathfrak{U}}}
%CONT>>
